% COSC 3337 Final Project -- Checkpoint 2: Database Design
% Due Sunday, November 8, 11:59 PM. Submit the compiled PDF on Canvas.
% Replace every [bracketed placeholder]. Put your ER diagram image in the same
% folder as this file and change the filename in Section 1.

\documentclass[11pt]{article}
\usepackage[margin=1in]{geometry}
\usepackage{graphicx}
\usepackage{enumitem}
\usepackage{tabularx}
\usepackage{booktabs}
\usepackage{amssymb}
\usepackage{xcolor}
\usepackage[colorlinks=true,urlcolor=blue!50!black]{hyperref}
\setlist{itemsep=2pt, topsep=4pt}

% ---------- Fill these in ----------
\newcommand{\appname}{[Application Name]}
\newcommand{\members}{[Member 1, Member 2, Member 3, Member 4]}
% -----------------------------------

% Primary keys are underlined; foreign keys are marked with \fk{...}.
\newcommand{\pk}[1]{\underline{#1}}
\newcommand{\fk}[1]{\textit{#1}$^{\mathrm{FK}}$}

\begin{document}

\begin{center}
  {\Large\bfseries COSC 3337 Checkpoint 2: Database Design}\\[4pt]
  {\large \appname{}}\\[2pt]
  \members
\end{center}

\section{ER Diagram}
% Use the ER notation from class. Include cardinality and participation
% constraints on every relationship.
% Replace the placeholder box with:
%   \includegraphics[width=\textwidth]{er_diagram.png}
\begin{center}
  \fbox{\parbox[c][3in][c]{0.9\textwidth}{\centering [Insert ER diagram image here]}}
\end{center}

\subsection*{Requirements check}
% Name the part of your diagram that satisfies each requirement.
\noindent\begin{tabularx}{\textwidth}{@{}lX@{}}
\toprule
Requirement & Where it appears in the diagram \\
\midrule
At least 6 entity sets & [List them] \\
At least one many-to-many relationship set & [Relationship name and the entities it connects] \\
At least one weak entity set & [Weak entity, its identifying relationship, and owner entity] \\
Cardinality \& participation on every relationship & [Yes / notes] \\
\bottomrule
\end{tabularx}

\section{Relational Schema}
% One line per table. Underline primary keys with \pk{} and mark foreign keys with \fk{}.
% Say what each foreign key references.
\begin{itemize}
  \item \textsc{[Table]}(\pk{[id]}, [attribute], [attribute], \fk{[other\_id]})
    \\ \quad \fk{[other\_id]} references \textsc{[OtherTable]}
  \item \textsc{[Table]}(\pk{[id]}, [attribute], [attribute])
  \item \textsc{[Table]}(\pk{[id]}, [attribute], [attribute])
\end{itemize}

\subsection*{How we translated the ER diagram}
\begin{itemize}
  \item \textbf{Weak entity:} [How the weak entity became a table, and what its primary key is]
  \item \textbf{Many-to-many relationship:} [The table you created for it, its primary key, and its foreign keys]
\end{itemize}

\section{Updated Questions}
% Your ten questions from Checkpoint 1. Revise them if the design changed, and
% note what changed. If nothing changed, say so.
\noindent\begin{tabularx}{\textwidth}{@{}rXl@{}}
\toprule
\# & Question & Changed? \\
\midrule
1  & [Question] & [No / how] \\
2  & [Question] & \\
3  & [Question] & \\
4  & [Question] & \\
5  & [Question] & \\
6  & [Question] & \\
7  & [Question] & \\
8  & [Question] & \\
9  & [Question] & \\
10 & [Question] & \\
\bottomrule
\end{tabularx}

\end{document}
